% Root and D full literal audits passed; independent rational control passed.
\subsection{Five states for the two-band four-element opposite corners}

\begin{lemma}[Unbounded rational cup vanishing for source order at most nine]
\label{C:nine-unbounded-act-cups}
For every finite monoid $M$ of order at most nine and every
finite right $M$-act, all rational degree-two products on its
ordinary act cone vanish. No bounded-homology assumption is
needed.
\end{lemma}
\begin{proof}
If $H^2(BM;\mathbb Q)=0$, the degree-two classes lift relative
to $BM$, and the suspension quotient kills their products.
Otherwise the minimal ideal has at least two rows and two
columns. At most two rows are covered by
Theorem~\ref{C:two-row-act-cups}. With at least three rows,
the order bound forces a trivial-group band of shape
$3\times2$ or $4\times2$, with at most three outside actors.
Lemma~\ref{C:three-outside-all-two-column-act-cups} applies
without a boundedness hypothesis and proves the assertion.
\end{proof}

\begin{lemma}[The rational five-state graph obstruction for small sources]
\label{C:nine-five-state-graph-square}
Let $M$ have order at most nine and act on a graph with vertex
set $U$ and oriented edge set $E$, where $|U|+|E|\leq5$.
Every actor sends each oriented edge at unit speed to another
oriented edge or collapses it to a vertex. If the cofiber of
its action-bar projection to $BM$ has one-dimensional second
rational cohomology, the square of that class is zero.
\end{lemma}
\begin{proof}
Regard the graph as the natural action graph for the right-zero
two-band on the state set $U\sqcup E$. The connected analysis
of Theorem~\ref{C:right-zero-five-connected} uses only
one-dimensional second rational cohomology, except for its
last literal control. That same thirteen-arrow control has
an exact rational square primitive as well.
The record
\path{code/outputs/a_rightzero_five_rational_control.json}
contains a nonboundary rational two-cocycle and a three-cochain
with only zero and one as coefficients. Its $872$ signed
four-simplex equations verify the square identity.
The standalone checker
\path{code/b_verify_rightzero_five_rational_control.py}
imports no producer or matrix-library routines. It rebuilds
all signed boundaries, checks their consecutive composites,
computes rational ranks $10,39$ for the first two
coboundaries, and verifies that the supplied cocycle is
nonboundary and satisfies all $872$ primitive equations.
Its complete record is
\path{code/outputs/b_rightzero_five_rational_control_verified.json}.
Thus the rational second cohomology is one-dimensional,
and the actual graph-image second-homology
isomorphism used in that proof also gives a rational
zero-square pullback. All connected cases are excluded
under the present rational hypothesis.

For a forest, the component projection is an equivariant
underlying equivalence, reducing to an ordinary act cone;
apply Lemma~\ref{C:nine-unbounded-act-cups}.
When every edge is a loop, the common vertex retraction gives
an ordinary-act homotopy retract and a suspension quotient.
If the retract has nonzero second cohomology its pullback
spans the assumed rank-one space and squares to zero. If it
has zero second cohomology, the class is relative and again
squares to zero. The remaining disconnected graphs are
exactly the three-vertex/two-edge cases in
Theorem~\ref{C:right-zero-five-state-reduction}. Its cycle
coinvariant, fixed-vertex, and two-eigenvalue arguments use
only the current rational rank-one assumption and the
two-state finite-line lemma. They therefore apply unchanged.
No intermediate space is asserted to have finite CW type.
\end{proof}

\begin{lemma}[An explicit normalized action-groupoid presentation]
\label{C:four-corner-normalized-action-groupoid}
Let $N=\{1,p,q,x\}$ have minimal right-zero ideal
$\{p,q\}$, and let $S$ be a left $N$-set. Write $P,Q,E$
for the actions of $p,q,x$, and let $U$ be the common image
of $P,Q$. In the fundamental groupoid of $N\ltimes S$,
identify each object $s$ with $Ps$ using the arrow labeled
$p$. The resulting groupoid on $U$ has one oriented
generator $g_s:Ps\to Qs$ for each $s\notin U$, with
$g_u=1$ for $u\in U$, and precisely these relations:
\[
\begin{array}{c|c}
\text{right action of $x$ on $\{p,q\}$}&\text{relations}\\\hline
\mathrm{id}&g_{Es}=g_s\\
\mathrm{constant}_p&g_{Es}=1\\
\mathrm{swap}&g_{Es}=g_s^{-1}.
\end{array}
\]
The constant case may be arranged to be $p$ by relabeling.
The presentation and its functor from the actual action
category are natural for every commuting right action.
\end{lemma}
\begin{proof}
All $N$-actors fix $U$ pointwise. In the fundamental groupoid,
all core endomorphisms vanish because $BN$ is contractible.
Normalize an arrow $n_s:s\to ns$ by the paths labeled $p$.
The equation $p_{ns}n_s=(pn)_s$ shows that its normalized
value is the identity if $pn=p$, and is $g_s$ if $pn=q$.
Thus every arrow is represented by the displayed generators.
The same equation with $q$ gives exactly the listed relation
for $E$; the relations for $P,Q$ only say $g_u=1$.

Conversely, define a functor from the action category to the
displayed groupoid by $s\mapsto Ps$, and send the arrow
$n_s$ to $1$ or $g_s$ according as $pn=p$ or $q$.
For $n=1$, use the identity. Each equation for a composable
pair follows from the displayed relation and the right
translation identity $\tau_{ab}=\tau_b\tau_a$.
In particular the $x^2$ equation adds nothing: in the first
two cases $x_s$ is sent to the identity; in the swap case
its product is $g_{Es}g_s=1$ and $x^2=1$.
This functor and the generator assignment into the completed
action category are inverse after the stated identifications
of objects. A commuting right map preserves $P,Q,E$ and
the equations, proving naturality.
\end{proof}

\begin{theorem}[Five states for the five two-band order-four corners]
\label{C:four-corner-two-band-five-states}
Suppose $|N|=4$, its minimal ideal is a right-zero band of
size two, $|M|\leq9$, and $|S|\leq5$ in an $(N,M)$-biset.
If the collage has one-dimensional second rational cohomology,
its generator has square zero. In particular none of these
collages has type $\mathbb CP^2$.
\end{theorem}
\begin{proof}
A singleton core $U$ is a common fixed state, so assume
$|U|\geq2$. Use the preceding presentation of the actual
left action category $Z=N\ltimes S$.

If $\tau_x$ is the identity, $E$ preserves both endpoints of
each graph edge. It is either an involution or a retraction.
The relations identify parallel edges in its orbits or with
their retraction images; an edge whose image is a core vertex
is a loop and is collapsed. The resulting groupoid is the
fundamental groupoid of an oriented graph $Q$ with the same
core vertices and at most $|S\setminus U|$ edges.
The right action sends an edge to an edge with the same
orientation or to a vertex, because it commutes with $E$.

If $\tau_x$ is constant, $PE=P$ and $QE=P$ show that each
edge indexed by an element of $ES\setminus U$ is a loop.
The relations simply collapse those loops, leaving another
oriented graph $Q$ with the same vertex and edge bounds.
The collapsed set is right-invariant, again giving a strict
right graph action. In these two cases the normalized
functor $Z\to\Pi Q$ induces isomorphisms on components and
first integral homology. The natural graph map
$Q\to B\Pi Q$ is an equivariant weak equivalence.
Consequently the actual collage maps to the $Q$ cofiber
through a natural equivalence zigzag, and its second
rational cohomology is the pullback of that cofiber's second
cohomology. The low-degree augmented bar comparison gives
an isomorphism here: the cone of the augmented chain map has
zero homology through degree two, since the target graph
has no second homology and the map is an isomorphism in
degrees zero and one. Its free-bar tensor has the same
connectivity. Lemma~\ref{C:nine-five-state-graph-square}
kills the square.

Finally suppose $\tau_x$ swaps $p,q$. Then $x^2=1$, so
$E$ is an involution fixing $U$ pointwise and reversing
the endpoints of each indexed edge. Each two-element
outside orbit contributes one unoriented edge. An outside
fixed point indexes a loop with relation $g_s^2=1$.
Kill these order-two loop generators. The resulting
groupoid is the fundamental groupoid of a graph $Q$ with
vertex set $U$ and at most
$\lfloor(|S|-|U|)/2\rfloor$ edges. This quotient is natural
for right actors: they preserve the set of $E$-fixed states,
and send a two-element edge orbit to another such orbit,
a fixed loop, or a vertex. Reversing an edge is allowed in
this last graph model.

The killed loops contribute only two-torsion to the
abelianization. Hence $Z\to\Pi Q$ still induces a
component bijection and a first rational homology
isomorphism, and the actual cofiber comparison induces
an isomorphism on second rational homology.
Since $|U|\geq2$ and $|S|\leq5$, this graph has at most
one edge. If that edge is a loop, it is the sole cycle.
An actor of degree other than $+1$ kills its rational
coinvariants, giving component-cone pullback and zero
products by Lemma~\ref{C:nine-unbounded-act-cups}.
Otherwise every actor fixes the unique loop vertex, giving
a common fixed state in the original collage. Thus the
loop case also has zero square.
For the forest case, the component projection reduces its
cofiber to an ordinary act cone. Apply
Lemma~\ref{C:nine-unbounded-act-cups}. Pullback along the
actual second-homology isomorphism finishes the proof.
\end{proof}

\begin{remark}[The remaining order-four shape]
\label{C:four-corner-five-state-partial-scope}
The theorem treats five of the six nonprojective right-trivial
four-element corners, and excludes their seventeen-arrow
profile $(8,4,5)$ as well as $(9,4,5)$ at eighteen arrows.
The identity-adjoined three-element right-zero band has a
different minimal ideal and is handled separately by
Theorem~\ref{C:right-zero-three-five-small-sources} below.
No actual finite CW type is inferred for any graph-quotient cofiber.
\end{remark}
